\customchapter{INTRODUCCIÓN}

% Las tablas de la Introducción se numeran con letras (Tabla A, B, ...)
% porque aquí el contador de capítulo todavía es cero.
\renewcommand{\thetable}{\Alph{table}}

Cuando se levanta una edificación patrimonial con un escáner láser terrestre, el
resultado es una nube de puntos $P=\{(x_i,f_i)\}_{i=1}^{N}$, donde
$x_i\in\mathbb{R}^3$ es la ubicación del punto y $f_i$ los atributos que registra
el escáner. Ya son la fuente primaria para el control geométrico, el modelado
como construido y la documentación de obra~\cite{wang2019construction}. Lo que no
aparece es la etiqueta. El archivo sabe dónde hay materia, pero no qué elemento
es cada cosa, y esa carencia bloquea el paso de la nube al modelo BIM y al
informe estructural o de conservación~\cite{wang2019construction}. Hoy ese
trabajo se hace a mano: construir el conjunto ArCH exigió anotar punto por punto
diecisiete escenas de edificios históricos~\cite{matrone2020benchmark}. No existe
un programa que automatice esa tarea sobre clases no catalogadas, y ese vacío es
el que este proyecto busca cubrir.

La nube de puntos guarda la geometría sin depender del ángulo ni de la
iluminación, a diferencia de las fotografías~\cite{pierdicca2020dgcnn}. Sobre esa
base, Pierdicca et al.~\cite{pierdicca2020dgcnn} adaptan una red de grafos
dinámicos con normales y color, con buenos resultados en muros y techos pero
caídas en columnas y arcos, y Matrone et al.~\cite{matrone2020benchmark}
confirman ese desbalance. Grilli y Remondino~\cite{grilli2020generalisation}
muestran el límite: un modelo entrenado sobre un edificio pierde desempeño en
otro de tipología distinta, porque sus descriptores geométricos no son estables
entre construcciones. Una tesis previa de UTEC~\cite{tesisutec2024} llega al
mismo resultado: su modelo dejó de reconocer columnas y arcos al pasar de un
patio a una capilla del mismo inmueble, y la única salida fue etiquetar a mano y
reentrenar.

Los métodos generales de segmentación semántica en tres dimensiones, como
PTv3~\cite{wu2024ptv3}, Sonata~\cite{wu2025sonata}, Stratified
Transformer~\cite{lai2022stratified} o Superpoint
Transformer~\cite{robert2023spt}, tampoco alcanzan. Trabajan con una lista de
clases fijada antes de entrenar y asignan todo punto imprevisto a la clase
conocida más parecida, sin avisar que no saben. En patrimonio eso es serio,
porque cada edificio tiene elementos propios que no figuran en ningún catálogo
previo~\cite{grilli2020generalisation}, y el problema se agrava con las clases
minoritarias~\cite{su2025brdl}.

El descubrimiento de clases novedosas, o NCD, elimina ese supuesto. Riz et
al.~\cite{riz2023nops} lo introdujeron para nubes de puntos y Xu et
al.~\cite{xu2024dasl} lo refinaron. Su mejor método actual es
HyperNCD~\cite{zhang2026hyperncd}, que resume la nube en \emph{prototipos},
vectores alrededor de los cuales se agrupan los puntos parecidos, y los conecta
mediante un \emph{hipergrafo}, donde cada arista une varios prototipos a la vez.
Rinde mal justo donde debería ser fuerte, con entre 11,8 y 47,5 mIoU en las
clases novedosas~\cite{zhang2026hyperncd}. Este trabajo señala como causa
probable el descriptor geométrico calculado a mano sobre el que arma sus
prototipos, el tipo de señal que Sonata~\cite{wu2025sonata} considera
perjudicial, y propone reemplazarlo por su codificador auto-supervisado
congelado. Además, el checkpoint entrenado con SemanticPOSS no se transfiere a la
catedral, pues concentra 98,4\,\% de los puntos en una sola cabeza; por ello M2
se entrenó con ArCH~\cite{matrone2020benchmark}, cuyo dominio corresponde al
patrimonio arquitectónico.

El Capítulo~I presenta la revisión crítica de la literatura; el Capítulo~II, el
marco teórico y el análisis de arquitectura; el Capítulo~III, el marco
metodológico; y el Capítulo~IV, los resultados obtenidos y sus límites.

\introsection{Formulación del problema}

Dado el levantamiento $P$ de una edificación patrimonial, con un subconjunto
etiquetado de clases conocidas $\mathcal{C}_k$ y otro sin etiquetar de clases
novedosas $\mathcal{C}_n$, donde $\mathcal{C}_k\cap\mathcal{C}_n=\emptyset$, el problema computacional consiste en producir una asignación
$\hat{y}:P\rightarrow\mathcal{C}_k\cup\mathcal{C}_n$ que clasifique bien los
puntos conocidos y descubra agrupaciones consistentes para los novedosos. El mejor método disponible
alcanza solo entre 11,8 y 47,5 mIoU sobre $\mathcal{C}_n$ y construye sus
prototipos sobre un descriptor geométrico analítico, con linealidad, planaridad y
dispersión, que la literatura auto-supervisada considera un atajo perjudicial. La
hipótesis de investigación es que reemplazarlo por una representación
auto-supervisada confiable aumenta de forma medible el mIoU sobre las clases
novedosas en un levantamiento patrimonial real.

\introsection{Objetivos de investigación}

\paragraph*{Objetivo General}
\addcontentsline{toc}{paragraph}{Objetivo General}

Construir y medir un sistema híbrido que integre las representaciones
auto-supervisadas de Sonata en el módulo de prototipos de HyperNCD, y cuantificar
la mejora del mIoU en la clasificación de elementos estructurales no anotados
sobre el levantamiento de la Catedral de Lima, mediante ramas congeladas y una
fusión posterior sin entrenamiento conjunto.

\paragraph*{Objetivos Específicos}
\addcontentsline{toc}{paragraph}{Objetivos Específicos}

El proyecto declara \textbf{dos objetivos con trabajo computacional propio}, OE1
y OE2, y \textbf{uno de salida analítica}, OE3.

\medskip
\noindent\textbf{OE1. Integrar Sonata en el módulo de prototipos.} Es el objetivo central. HyperNCD describe cada punto con un vector $F_i$ que
junta $f_{\mathrm{sem}}$, el descriptor semántico de su red supervisada, y
$f_{\mathrm{geo}}$, el geométrico calculado a mano (Ec.~\eqref{eq:hyperncd3}). Se
implementa el bloque de fusión que sustituye $f_{\mathrm{geo}}$ por
$f_{\mathrm{ssl}}$, el descriptor del codificador congelado de Sonata, y resuelve
la diferencia de resolución entre PTv3 y la rejilla de Minkowski con
\emph{up-cast} y \emph{cache} por sección.

\begin{itemize}\setlength{\itemsep}{0pt}
\item \textit{Herramientas y trabajo computacional:} pesos públicos de Sonata,
Pointcept, PTv3, \texttt{spconv}, MinkowskiEngine y PyTorch; módulo de fusión e
inferencia congelada sobre cada sección.
\item \textit{Resultado verificable:} repositorio con ambas ramas congeladas,
archivos de descriptores y el híbrido produciendo
$F_i=[f_{\mathrm{sem}};f_{\mathrm{ssl}}]$ sobre una sección de control.
\item \textit{Conclusión asociada:} si la sustitución es viable dentro del
hipergrafo.
\end{itemize}

\medskip
\noindent\textbf{OE2. Evaluar la configuración del híbrido.} Comparar M2, M3,
M4 base y M4 bajo un protocolo común, cuantificando la variabilidad debida a las
semillas de agrupamiento. Los parámetros activos se documentan desde el código
y los registros de ejecución.

\begin{itemize}\setlength{\itemsep}{0pt}
\item \textit{Herramientas y trabajo computacional:} entrenamiento de M2 con
ArCH, inferencia congelada, agrupamiento con diez semillas y métricas mIoU, NMI
y dominancia.
\item \textit{Resultado verificable:} tabla comparativa de las cuatro
configuraciones, con media y desviación estándar del agrupamiento.
\item \textit{Conclusión asociada:} si la fusión mejora numéricamente la línea
base y si la diferencia supera la dispersión observada.
\end{itemize}

\medskip
\noindent\textbf{OE3. Cuantificar la mejora y delimitar el rango de operación.}
\textbf{No incorpora trabajo computacional nuevo}. Es la salida de análisis de OE1
y OE2. Consiste en medir el $\Delta$mIoU de las clases novedosas
(Ec.~\eqref{eq:delta}) frente a la línea base, identificar qué variable domina
(Ec.~\eqref{eq:dominante}), comparar las tres variantes del descriptor
, geométrica, auto-supervisada y concatenación,  y fijar el rango favorable.

\begin{itemize}\setlength{\itemsep}{0pt}
\item \textit{Herramientas:} el mIoU con emparejamiento húngaro, que decide qué
grupo descubierto corresponde a qué clase real; los índices ARI y NMI; la brecha
cerrada $G$ (Ec.~\eqref{eq:G}); y visualizaciones PCA.
\item \textit{Resultado verificable:} tabla comparativa, estudio de ablación y
curva de degradación que fija el rango declarado.
\end{itemize}

\introsection{Justificación}

La clasificación de elementos estructurales se hace hoy a mano porque no hay
herramienta que cubra la tarea cuando la lista de clases no está cerrada, y cada
punto de mIoU ganado es corrección que el especialista deja de hacer. Los
usuarios están conformes con la precisión que obtienen a mano, así que el valor
está en automatizar la tarea a un nivel suficiente, no en igualar esa precisión.

El aporte no es una arquitectura nueva. Lo que se busca es establecer si un problema ya
demostrado con supervisión completa se mantiene cuando no hay etiquetas que
corrijan al modelo. Es un resultado incierto, medible y que ningún trabajo previo
ha reportado, y como se mueve una sola pieza, cualquier variación se puede
atribuir al reemplazo.

\introsection{Alcance y limitaciones}

\paragraph*{Objeto de estudio y generalización}
Se trabaja sobre un único edificio, la \textbf{Catedral de Lima}, con el
levantamiento provisto por la Facultad de Ingeniería Civil. El objeto es fijo y
el modelo se ajusta a él, no al revés. Los resultados se declaran válidos para tipología de catedral o iglesia colonial,
es decir nave con columnas y arcos, cubierta abovedada y fachada con torres, y no
se extienden a vivienda ni a obra vial o industrial.

\paragraph*{Volumen de datos}
El levantamiento se recibe como archivos LAS (Sección~\ref{sec:las}) y la unidad
de trabajo es la \textbf{sección espacial} (Sección~\ref{sec:seccion}). Una sección es el conjunto de
archivos que se preprocesan juntos, y las secciones funcionan como particiones de
entrenamiento y evaluación. El preprocesamiento
aplica dos reducciones de naturaleza distinta: el filtrado SOR descarta los
puntos estadísticamente aislados y la voxelización agrupa los puntos de cada cubo
de 5\,cm en un solo representante, sin descartar ruido. Las cifras medidas están
en la Tabla~\ref{tab:m1-metadatos}. El vóxel de 5\,cm es un límite, no una
preferencia: por debajo, con muy pocos puntos por clase novedosa, la diferencia
entre el modelo base y el híbrido dejaría de distinguirse del ruido entre
corridas.

\paragraph*{Hardware y costo}
El proyecto está limitado por cómputo y \textbf{no estima tiempos de calendario}.
El hardware en el que las pruebas corrieron realmente y los tiempos y memorias
medidos están en la Tabla~\ref{tab:res-hardware} y en el
\hyperref[sec:anexo-software]{Anexo~A}. La división en secciones sostiene esta
política, ya que ninguna etapa necesita la nube completa en memoria.

\paragraph*{Rangos de operación}
El proyecto declara por adelantado dónde se espera que el método funcione, y OE3
lo verifica. Se espera que funcione con densidad de al menos 1\,pt/cm$^2$, con
elementos de 20\,cm o más de firma geométrica repetida, como columnas, arcos,
bóvedas, cornisas y contrafuertes, y hasta seis clases novedosas por sección. No
se espera buen comportamiento con ornamento fino de menos de 5\,cm, superficies
muy ocluidas, ni clases novedosas bajo el 0,5\% de los puntos, régimen en el que
HyperNCD colapsa a 11,8 mIoU~\cite{zhang2026hyperncd}, ni con muchas clases simultáneas,
donde Sonata baja a 29,3 mIoU en ScanNet200 frente a 72,5 en
ScanNet~\cite{wu2025sonata}.

\paragraph*{Exclusiones}
Quedan fuera el preentrenamiento auto-supervisado propio, que en Sonata exigió
140 mil escenas y 32 GPU, la detección de instancias, el vocabulario abierto y la
generación del modelo BIM.

% Restauración de la numeración de tablas por capítulo.
\renewcommand{\thetable}{\arabic{chapter}.\arabic{table}}
\setcounter{table}{0}
